% Explanatory notes 1-38, from john-dowland-letter-modern-english.docx.

\note{1}{Right Honourable and bounden duty}{Dowland addresses Robert Cecil, then a Privy Councillor. ‘Bounden’ means morally and personally obliged or indebted, not physically restrained. The opening combines conventional humility with a practical appeal to Cecil as patron and political protector.}

\note{2}{Preserve and confound}{To ask for God’s ‘preservation’ is to ask that God protect and sustain the Queen and country. ‘Confound’ here means frustrate, defeat or overthrow, rather than merely confuse. The pronoun ‘them’ refers jointly to Queen and country.}

\note{3}{Sir Henry Cobham and lay in Paris}{Cobham was the English ambassador in Paris from 1579 to 1583. Dowland’s being his ‘servant’ means that he belonged to the ambassador’s household or service; it need not mean that he performed ordinary domestic labour. To ‘lie’ in a city could mean to reside or be stationed there.}

\note{4}{The Paris acquaintances}{Nicholas Morgan had formerly been a gentleman of the Chapel Royal, an institution of royal musicians and clergy rather than simply a physical chapel. Richard Verstegan was a Catholic writer and exile. Morris was the household porter or doorkeeper; ‘at Rome’ was a normal idiom for residing in Rome.}

\note{5}{Idle toys and overreached}{‘Idle toys’ means foolish ideas, fantasies or trifles; ‘toys’ could mean notions, not only children’s objects. Dowland’s unusual spelling ‘overecht’ is best understood as ‘overreached’: they outwitted, misled or deceived him. ‘I believed with them’ means that he adopted or shared their beliefs.}

\note{6}{Within two years after}{The phrase means that within two years after his Paris experience Dowland came back to England, not that the executions occurred within two years after his arrival. ‘That faction’ means the Catholic party or group. Dowland says that at the time he believed religion alone caused the punishments; the letter later presents this belief as youthful error.}

\note{7}{Mr Johnson’s place}{John Johnson, one of the leading English lutenists, served in the royal musical establishment and died in 1594. Dowland became a ‘suitor’ for his place: an applicant or petitioner. His aside that he thought himself ‘most worthiest’ reveals both confidence in his ability and resentment at the outcome.}

\note{8}{Preferment and sounding the cause}{‘Preferment’ means appointment, promotion or advancement through patronage. To ‘sound’ the cause means to probe or investigate it. Dowland says he guessed that his reputation as a Catholic prevented his appointment; he does not prove that this was the official reason.}

\note{9}{Beyond the seas and the Privy Council}{‘Beyond the seas’ means abroad, especially continental Europe. ‘Durst’ is the old past tense of ‘dare’. The Privy Council was the Queen’s principal executive and advisory body; ‘privy’ means private or confidential. Unlicensed foreign travel could attract suspicion, particularly amid fears of Catholic conspiracy.}

\note{10}{Both your hands}{The phrase most likely means that Cecil and the Earl of Essex supplied their signatures, warrants or written approval. It does not mean that they physically gave Dowland both hands. His parenthesis offers future service in repayment.}

\note{11}{The Duke of Brunswick’s gifts}{The duke was Heinrich Julius of Brunswick-Luneburg, whose court was at Wolfenbuttel. The chain, money and luxury textiles were marks of princely favour and potential employment. ‘Apparel’ means clothing; historically the word is related to equipping or preparing a person.}

\note{12}{One of my quality}{Here ‘quality’ means rank, profession, accomplishment or social standing. Dowland is saying that the welcome was exceptionally great for a musician of his position, not commenting on the abstract quality of his music.}

\note{13}{The Landgrave’s gifts}{Moritz, Landgrave of Hesse-Kassel, was himself musically accomplished. A ‘standing cup’ was a large ceremonial drinking vessel with a foot and often a lid. The ‘dollars’ were thalers. The gifts to Dowland and his wife demonstrate the status he enjoyed abroad.}

\note{14}{The Duke in Florence}{This was Ferdinando I de’ Medici, Grand Duke of Tuscany. ‘Got great favours’ means that Dowland received a favourable reception and patronage, not necessarily one particular formal reward.}

\note{15}{Skidmore or Scudamore}{The priest was John Scudamore, whose name Dowland spells in several ways. He was an English Catholic priest, not an Anglican clergyman. The unnamed gentleman in the piazza merely pointed him out; the gentleman plays no further part in the story.}

\note{16}{Luca Marenzio}{Marenzio was one of the leading Italian madrigal composers of the period and was associated with Rome. Dowland says that musical study was the purpose of his intended journey. The letter does not say that Skidmore would travel with him, and Dowland apparently did not reach Rome on this journey.}

\note{17}{Mistaken in Rome}{Dowland probably fears being mistaken for a Protestant agent, spy or religious enemy in Catholic Rome. The wording is deliberately general. He presents himself as caught between two suspicions: treated as Catholic in England but potentially mistrusted in Rome.}

\note{18}{Obstinate papist}{‘Papist’ was a hostile Protestant term for a Roman Catholic. The Queen’s alleged statement is hearsay reported by Dowland, not a surviving direct quotation from Elizabeth. It praises his musical usefulness while identifying his religion as a political obstacle.}

\note{19}{Make her words true}{Skidmore’s reply means: if you are not yet an obstinate Catholic, become one and thereby make the Queen’s description accurate. It is an invitation to commit himself more fully to Catholicism.}

\note{20}{Touch her finger and meddle nor make}{To say that priests would not even ‘touch her finger’ intensifies Dowland’s disbelief that they would harm Elizabeth. ‘Neither meddle nor make with’ means neither become involved nor take part. His qualification is political: he will stay uninvolved so long as nothing is done against the Queen.}

\note{21}{Spanish faction jars and the pike}{The ‘Spanish faction’ means Catholics aligned politically with Philip II of Spain. ‘Jars’ are quarrels or conflicts. To ‘bear a pike’ means to serve as an ordinary infantryman. Skidmore presents himself as Catholic but English and anti-Spanish, illustrating divisions among exiled English Catholics.}

\note{22}{Order to attach}{In sixteenth-century usage ‘attach’ commonly meant arrest, seize or take into legal custody. That is the likeliest sense, although Dowland gives too little context to establish what authority Skidmore claimed or which gentlemen were meant.}

\note{23}{Keep in and write letters of me}{To ‘keep in’ probably means remain indoors, secluded or confined for a month. ‘Write letters of me’ can mean write about Dowland or write recommendations on his behalf. Santa Maria Novella was a major Dominican religious house in Florence.}

\note{24}{Friar Baylie}{A friar is a member of a Catholic mendicant religious order; the closest German terms are ‘Ordensbruder’ or, more specifically, ‘Bettelmönch’. It does not simply mean any Catholic priest, but Baylie is certainly presented as a Catholic religious.}

\note{25}{Warning Lord Grey}{Dowland says he immediately warned English gentlemen about Skidmore’s information, whereupon they separated. This episode helps him present himself to Cecil as already acting loyally before he wrote the letter.}

\note{26}{Josias Bodley as witness}{Dowland stresses that he opened Baylie’s letter before Josias Bodley and other witnesses. He also enclosed the letter for Cecil. These details provide corroboration and make the account look like an intelligence report rather than secret Catholic correspondence.}

\note{27}{The papal offer}{Baylie’s promises address precisely Dowland’s grievances: lack of employment, insufficient money and exclusion from honour. Whether a pension had truly been authorized is unknown; within Dowland’s narrative, the offer represents the height of the temptation.}

\note{28}{Who might lose their lives}{The pronoun is ambiguous. It most probably refers to the agents who entered England to escort Catholics out and who risked execution if caught. It could grammatically be taken to refer to Dowland’s family, but the following sentence about agents entering England favours the first reading.}

\note{29}{Seventeen people and the bishop}{This is Baylie’s claim, repeated by Dowland. The bishop is not identified. The story praises Skidmore for helping English Catholics escape, but Dowland now offers it to Cecil as evidence of clandestine activity.}

\note{30}{For want}{Here ‘want’ means poverty, need or lack of livelihood. Dowland says that professional exclusion might drive him into papal service. The passage is also carefully self-defensive: material desperation, rather than disloyal intention, has brought him close to England’s enemies.}

\note{31}{Treason treachery and Mass}{‘Treason’ is betrayal of the sovereign or state and had a specific legal force; ‘treachery’ is the broader idea of deceitful betrayal. The grammar most naturally makes Mass the ‘great abuse of the people’, perhaps because its Latin ritual was unintelligible to ordinary worshippers and, Dowland says, even to him. The sentence is polemical and may also be a calculated profession of Protestant sympathy.}

\note{32}{Reformed myself}{In context this certainly promises obedience to Elizabeth’s religious laws. Whether it proves an inward theological conversion is less certain. Dowland may be promising outward conformity and lawful conduct rather than describing a complete change of private belief.}

\note{33}{Pierce and Dracot}{Dowland gives no accusation or further story about these men. Their names may simply be intelligence leads for Cecil’s network. Their mention is relevant because both were travelling to Rome, but their precise role remains unknown.}

\note{34}{King Philip}{This is Philip II of Spain, Catholic king and former husband of Mary I of England. Spain and England were at war; the Spanish Armada had attempted an invasion in 1588, and further Spanish operations remained a serious fear. ‘Dispatch’ here means kill.}

\note{35}{Forsaking them and their religion}{This sounds like an explicit renunciation of Catholicism. Nevertheless, the letter as a whole is a plea to a powerful Protestant statesman, and its sincerity cannot be measured directly. The safer conclusion is that Dowland publicly rejects the Catholic network described here and promises conformity and loyalty to Elizabeth.}

\note{36}{Skidmore’s whereabouts}{Dowland infers that Skidmore and the priests may be in England because Skidmore did not remain in Florence as promised. Baylie’s statement that Skidmore went to France to study law is reported but apparently not trusted. Dowland is offering suspicion, not confirmed knowledge.}

\note{37}{Instructions through my wife}{Dowland asks Cecil to communicate with him through Dowland’s wife in England. His offer to ‘lose my life against them’ means to fight, and if necessary die, against a Spanish attack. It is chiefly a pledge of loyalty, not a clearly formulated request to enlist in a particular military unit.}

\note{38}{Ill writing and worse inditing}{‘Writing’ refers to handwriting; ‘inditing’ means composing or wording the letter. The conventional apology closes a broader request: Cecil should pardon Dowland’s presumption, accept his account, restore his reputation and consider him available for English service.}
